\section{Conclusion}\label{sec:conclusion}

The Bailout Protocol, as specified in this paper, establishes the mandatory architectural mechanism by which the \bxthree{} Framework makes its upstream accountability guarantee structurally operative at runtime. The protocol is not an error-handling feature selected by a system at its discretion; it is an architectural compulsion embedded in the \fact{} layer of every \bxthree{} node, enforced by the deterministic transition relation of the state machine $\mathcal{B}$, and rendered irrevocable by the parent-chain bypass that prohibits any machine actor from absorbing, resolving, or suppressing an escalated exception. The human anchor $h(n)$ is not a design choice within this architecture --- it is a structural necessity imposed by the functional separation between the \purpose{}, \bounds{}, and \fact{} layers, and the protocol is the mechanism by which this structural necessity is enforced at runtime.

The protocol's three correctness properties --- Safety, Liveness, and Accountability --- are not aspirational targets but formal consequences of the specification. Safety follows from the atomic and irreversible nature of the \texttt{ESCALATING} transition and the bypass mechanism's prohibition on machine-actor absorption. Liveness follows from the timeout architecture's guarantee of bounded progression through each protocol phase and the \fact{} layer's enforcement of the state machine's transition relation. Accountability follows from the immutable parent pointer's guarantee of a structurally fixed escalation path and the Forensic Ledger's append-only record of every protocol event and resolution. Together, these properties define the protocol's correctness contract: every unanticipated exception that exceeds a node's operating range $R(n)$ and cannot be resolved within $T_{\text{bounds}}$ will, with mathematical certainty, reach a named human principal who is identifiable at the protocol's terminus, and the resolution will be recorded as an immutable architectural record.

The contribution of the Bailout Protocol to the \bxthree{} Framework is specific and irreplaceable. The framework's three-layer architecture defines the functional separation between intent, bounded reasoning, and deterministic enforcement, but the separation alone does not guarantee that exceptions propagate to a human principal rather than being absorbed at an intermediate machine tier. Without the Bailout Protocol, the framework's accountability guarantee would be a claim without a mechanism --- structurally sound but operationally unenforceable. With the Bailout Protocol, the guarantee is enforceable for every unanticipated exception condition across the full recursive depth of the agent population, and the enforcement is verifiable through the Forensic Ledger by any auditor with access to the \fact{} layer's records.

The Bailout Protocol distinguishes the \bxthree{} Framework from prior escalation architectures by making human oversight compulsory rather than optional, and by enforcing this compulsion through architectural constraint rather than through behavioral convention. The distinction is not incremental. A system that may consult a human is not equivalent to a system that cannot proceed without one; the former permits accountable absorption at intermediate tiers, while the latter establishes human contact as an architectural precondition for the continuation of autonomous operation. This distinction is the core contribution of the Bailout Protocol, and it is the property that makes the \bxthree{} Framework's upstream accountability guarantee structurally operative rather than merely claimed.

The limitations identified in Section~\ref{sec:limitations} --- human response latency, adversarial exploitation risk, scalability of the human anchor, and escalation pathway weaponization --- are real and operationally significant. They do not undermine the protocol's correctness properties as stated under the specified assumptions; rather, they define the boundary of the protocol's applicable scope and motivate the directed research agenda described in the future work section. Formal verification, adaptive timeout mechanisms, and distributed human anchor pools are not speculative extensions but necessary steps toward a protocol that is both theoretically sound and operationally robust in real-world deployments at scale. The Bailout Protocol provides the foundation. The research roadmap provides the path.